% Add after the global n=2 two-zero theorem in lerch_boundary_notes.tex.
% No numerical input beyond the exact endpoint sign already certified there.

\subsection{Monotonicity of the larger branch and a unique level crossing}
\label{lerchshape:sec:n2}

The single change of coefficient sign also determines the shape of the
larger zero branch and the part of the smaller branch above $a=1$.
These conclusions leave a sharper question about the number of extrema
of the smaller branch below this level.

\begin{corollary}[Shape constraints for the two zero branches]
\label{lerchshape:cor:n2}
Let $a_1(\rho)<a_2(\rho)$ be the two simple positive zeros in
Theorem~\ref{lerchboundary:thm:n2global}. Then $a_2$ is strictly
decreasing on $[0,1]$, with
\[
 \frac{13}{10}<a_2(\rho)<e^2\qquad(0<\rho\leq1),
 \qquad a_2(0)=e^2.
\]
There is a unique $\rho_c\in(0,1)$ such that $a_1(\rho_c)=1$.
Moreover,
\[
 a_1(\rho)<1\quad(0<\rho<\rho_c),\qquad
 a_1(\rho)>1\quad(\rho_c<\rho\leq1),
\]
and
\[
 a_1'(\rho)>0\quad(\rho_c\leq\rho<1),\qquad
 a_2'(\rho)<0\quad(0<\rho<1).
\]
In particular, the level crossing is transverse, and every interior
critical point of $a_1$ lies in $(0,\rho_c)$.
\end{corollary}

\begin{proof}
Write $G=G_{2,1}$ and
\[
 f(u)=\frac{\log u(\log u-2)}{u^2},\qquad
 G(a,\rho)=\sum_{m=0}^{\infty}\rho^m f(a+m).
\]
For a fixed $a\in[1,e^2)$ set $c=e^2-a>0$ and $b_m=f(a+m)$.
The coefficient $b_m$ is nonpositive for $m<c$ and nonnegative
for $m>c$, and it vanishes if $m=c$. Thus
\[
 (m-c)b_m\geq0\qquad(m\geq0).
\]
There are strictly positive terms, for example every integer $m>c$.
Termwise differentiation is valid for $0<\rho<1$, and gives the
strict differential inequality
\begin{equation}\label{lerchshape:eq:weighted}
 \rho G_\rho(a,\rho)-cG(a,\rho)
   =\sum_{m=0}^{\infty}(m-c)b_m\rho^m>0.
\end{equation}
The use of a real cutoff $c$ includes without exception the possibility
that one coefficient vanishes at $a+m=e^2$.

For $\rho>0$, the $m=0$ term at $a=e^2$ vanishes and all terms with
$m\geq1$ are positive. Therefore $G(e^2,\rho)>0$. Together with
$G(13/10,\rho)<0$ and the two-zero theorem, this places the larger
zero in $(13/10,e^2)$, including at $\rho=1$. The signs at the
simple zeros are
\[
 G_a(a_1(\rho),\rho)<0,
 \qquad G_a(a_2(\rho),\rho)>0,
\]
since $G$ is positive before the smaller zero, negative between
the two zeros, and positive after the larger zero. At any zero
in $[1,e^2)$, \eqref{lerchshape:eq:weighted} gives $G_\rho>0$.
Implicit differentiation consequently gives $a_2'=-G_\rho/G_a<0$
for $0<\rho<1$. Continuity at both endpoints proves strict decrease
on the whole closed interval.

Now fix $a=1$ and put $c_0=e^2-1$. The same inequality shows that
\[
 \frac{d}{d\rho}\bigl(\rho^{-c_0}G(1,\rho)\bigr)
 =\rho^{-c_0-1}\bigl(\rho G_\rho(1,\rho)-c_0G(1,\rho)\bigr)>0
 \qquad(0<\rho<1).
\]
The expression is continuous up to $\rho=1$. Since $f(1)=0$ and
$f(2)<0$, one has $G(1,\rho)=\rho f(2)+O(\rho^2)<0$ for small
positive $\rho$. The exact certificate
\eqref{lerchboundary:eq:n2cert-one} gives $G(1,1)>0$.
Strict increase of the weighted expression therefore gives precisely
one zero $\rho_c\in(0,1)$, with negative sign before and positive
sign after it.

Because $1<13/10<a_2(\rho)$, the sign of $G(1,\rho)$ is negative
exactly when $a_1(\rho)<1$, and positive exactly when $a_1(\rho)>1$.
This proves the asserted unique crossing and the inequalities on
each side. When $\rho_c\leq\rho<1$, the smaller root lies in
$[1,13/10)$, so \eqref{lerchshape:eq:weighted} applies there as well.
Its negative $a$ derivative now gives $a_1'=-G_\rho/G_a>0$,
including at the crossing itself.
\end{proof}

The remaining shape question is whether $a_1$ has exactly one minimum
in $(0,\rho_c)$; a sharper version asks whether its derivative has
exactly one zero there. Neither uniqueness assertion is required for
the corollary.
